\rok{Новембар 2023.}{F}

Први практични тест из Алгоритама и структура података

\zadatak{1}{Додавање елемента на крај циркуларне листе}
Написати методу која имплементира стандардан начин за додавање елемената на крај циркуларне листе.

\textbf{Додатне напомене за решавање задатка:}
\begin{itemize}
    \item Алгоритам написати у оквиру закоментарисане методе \texttt{addTail(int value)}.
    \item Поменута метода се налази у оквиру класе \texttt{CircularList} која се налази у придруженом шаблону.
    \item Обезбедити код у \texttt{Main} методи за тестирање алгоритма.
\end{itemize}



\zadatak{2}{Инвертовање низа коришћењем структуре стека}
Креирати функцију која ће статички низ бројева инвертовати коришћењем структуре стека (класа \texttt{Stack}). Алгоритам је неопходно написати са линеарном временском комплексношћу $O(n)$.

\textbf{Додатни захтеви за решавање задатка:}
\begin{itemize}
    \item Алгоритам не сме креирати нови низ већ инвертовати елементе постојећег.
    \item Захтеван потпис методе:
    \begin{Verbatim}[fontsize=\footnotesize, numbers=left, xleftmargin=10mm]
public static void InvertArrayUsingStack (long[] arr)
    \end{Verbatim}
\end{itemize}

\textbf{Примери:}
\begin{itemize}
    \item \textbf{Улаз:} \texttt{[2,3,5,4,10]} \\
          \textbf{Излаз:} \texttt{[10,4,5,3,2]}
    \item \textbf{Улаз:} \texttt{[]} \\
          \textbf{Излаз:} \texttt{[]}
\end{itemize}



\zadatak{3}{Кришнамуртијев број}
Креирати алгоритам који проверава да ли се од елемената циркуларне листе може креирати Кришнамуртијев број. Кришнамуртијев број је број чија је сума факторијела сваке цифре једнака самом броју. Алгоритам је неопходно написати са линеарном временском комплексношћу $O(n)$.

\textbf{Додатни захтеви за решавање задатка:}
\begin{itemize}
    \item Неопходно је имплементирати функцију \texttt{Factorial} која треба да пронађе факторијел прослеђеног броја.
    \item Број 145 је Кришнамуртијев број: 1! + 4! + 5! = 1 + 24 + 120 = 145
    \item Број 15 није Кришнамуртијев број: 1! + 5! = 1 + 120 = 121
    \item Задатак решавати помоћу структуре једноструко спрегнуте листе (класа \texttt{LinkedList}) која је дата у оквиру шаблона задатка. Алгоритам је неопходно написати са линеарном временском комплексношћу $O(n)$.
    \item Захтеван потпис методе:
    \begin{Verbatim}[fontsize=\footnotesize, numbers=left, xleftmargin=10mm]
public void IsListKrishnamurthyNumber()
    \end{Verbatim}
\end{itemize}

\textbf{Примери:}
\begin{itemize}
    \item \textbf{Улаз 1:} $1 \rightarrow 4 \rightarrow 5$ \\
          \textbf{Излаз 1:} \texttt{"Number 145 is Krishnamurthy number."}
    \item \textbf{Улаз 2:} $1 \rightarrow 2 \rightarrow 5$ \\
          \textbf{Излаз 2:} \texttt{"Number 125 is Krishnamurthy number."}
\end{itemize}

\sablon{Шаблон првог теста новембар 2023. група F}{2023/Novembar/GrupaF/AISP2023_Test1_GrupaF.linq}
